\section{Analysis}\label{sec:analysis}

Throughout, $\Theta_{g,k}=V w_g\eta_k$ and $B_{g,k}=\Theta_{g,k}+\rmax_k$, with $\eta_k$ from \eqref{eq:eta}.
We assume that the per-slot problem is solved exactly (Proposition~\ref{prop:solver}).

\subsection{A deterministic bound on class envy}\label{sec:envy}

\begin{theorem}[Deficit bound]\label{thm:bound}
For every sequence of capacities, arrivals, departures and demand caps, and for arbitrary real tier prices $\beta_{g,t}$, \repay{} satisfies
\[
D_{i,k}(t)\le B_{g_i,k}=V w_{g_i}\eta_k+\rmax_k
\]
for every elastic device $i$, every class $k\in\cK$ and every slot $t$ of its sojourn.
\end{theorem}

\begin{proof}
Fix $i$, $k$, a slot $t$, and write $g=g_i$, $D=D_{i,k}(t)$, $\lambda$ for the optimal capacity multiplier.
We first show that $D\ge\Theta_{g,k}$ implies $r_{i,k,t}\le x_{i,t}$.

If the group $(g,k)$ is empty or $a_{g,k}=0$, then $r_{i,k,t}=0\le x_{i,t}$.
If $x_{i,t}=d_{i,t}$, then $r_{i,k,t}=a\min(z,d_{i,t})\le d_{i,t}=x_{i,t}$.
Suppose then that $a_{g,k}>0$ and $x_{i,t}<d_{i,t}$.
Since $x_i$ can be increased, stationarity of the Lagrangian in $x_i$ gives $Vw_gu'(x_i)+\beta_g+D_i-\lambda\le0$, and because $u'>0$ and $D_i\ge D$,
\begin{equation}\label{eq:pf1}
\lambda-\beta_g>D .
\end{equation}
The Lagrangian term of the group, $n\,Vw_g\,a\,u_k(y/a)+(n\beta_g-\pi-n\lambda)y-\zeta a$, is positively homogeneous in $(a,y)$ and is maximised over a cone intersected with $a\le1$; a maximiser with $a>0$ must therefore have a nonnegative value, that is, with $z=y/a$,
\[
n\,Vw_g\,u_k(z)\ \ge\ n(\lambda-\beta_g)z+\pi z+\zeta\ \ge\ n(\lambda-\beta_g)z ,
\]
where we used $\pi,\zeta\ge0$. Dividing by $nz>0$,
\begin{equation}\label{eq:pf2}
\lambda-\beta_g\le Vw_g\frac{u_k(z)}{z}\le\Theta_{g,k}.
\end{equation}
Combining \eqref{eq:pf1} and \eqref{eq:pf2} gives $D<\Theta_{g,k}$.
So if $D\ge\Theta_{g,k}$ one of the first two cases applies and $r_{i,k,t}\le x_{i,t}$.
Rounding does not alter this: it never lowers an elastic allocation, and the reference in \eqref{eq:D} is the fluid one.

Now induct on $t$. The deficit starts at $0$. If $D_{i,k}(t)<\Theta_{g,k}$, then $D_{i,k}(t+1)\le D_{i,k}(t)+r_{i,k,t}<\Theta_{g,k}+\rmax_k$. If $\Theta_{g,k}\le D_{i,k}(t)\le B_{g,k}$, then $D_{i,k}(t+1)\le D_{i,k}(t)\le B_{g,k}$.
\end{proof}

Three features of the proof are worth pointing out.
First, it uses no assumption on the inputs: the bound holds on adversarial sequences and for populations that change arbitrarily.
Second, the tier price $\beta_g$ cancels because it multiplies the rate of elastic and RT flows of the tier alike; this is why we wrote the business constraints per unit of rate. A constraint that rewarded only RT admissions of a tier would add its multiplier to the bound.
Third, only the elastic bonus is used. The debt price on RT groups is not needed for the envy bound; it is what makes the rule a true drift-plus-penalty method and hence near-optimal (Section~\ref{sec:opt}).

\begin{remark}[A sharper constant]\label{rem:sharp}
If the device is not rate-limited below the RT rate ($d_{i,t}\ge\rmax_k$, so $s_{i,k}=1$), then $\pi_{g,k}\ge D$ and the chain above gives $n\,Vw_gu_k(z)>(n+1)Dz$, hence growth stops already at $D\ge\Theta_{g,k}\,n_{g,k}/(n_{g,k}+1)$.
\end{remark}

\begin{corollary}[Window class-envy-freeness]\label{cor:window}
On every sample path, for every elastic device $i$, class $k$ and slots $s\le t$ in its sojourn,
\[
\sum_{\tau=s}^{t}x_{i,\tau}\ \ge\ \sum_{\tau=s}^{t}r_{i,k,\tau}-\big(D_{i,k}(t+1)-D_{i,k}(s)\big)\ \ge\ \sum_{\tau=s}^{t}r_{i,k,\tau}-B_{g_i,k}.
\]
In particular \repay{} is $B$-class-envy-free with $B=\max_{g,k}B_{g,k}$, and over a sojourn of $L$ slots the envy per slot is at most $B/L$.
\end{corollary}
\begin{proof}
By \eqref{eq:D}, $D_{i,k}(\tau+1)\ge D_{i,k}(\tau)+r_{i,k,\tau}-x_{i,\tau}$. Sum over $\tau=s,\dots,t$ and apply Theorem~\ref{thm:bound}.
\end{proof}

The bound is a constant in units of volume. With $V=20$, $w_g=1$ and $\rmin_V=0.5$~Mbit/s it equals $43$~Mbit, about fourteen seconds of a full-rate video call, whatever the length of the sojourn.

\subsection{Incentives}\label{sec:ic}

Envy-freeness compares what a device received with what the other class was offered \emph{on the realised path}.
A misreport would change the path: the deviating device joins an RT group, and its own bonus and debt terms disappear.
To pass from envy to incentives we follow the large-market approach of \citet{azevedo2019strategy}, who show that envy-freeness is a sufficient condition for a mechanism to be strategy-proof in the large.

\begin{definition}[Price-taking]\label{def:pt}
A device is \emph{price-taking} if it evaluates a deviation assuming that the multipliers and offers $(\lambda_t,\beta_{g,t},a_{g,k,t},z_{g,k,t})_t$ of the truthful path stay unchanged, and that in each slot in which it declares class $k$ it receives the offer of group $(g,k)$, that is, an expected usable volume of at most $r_{i,k,t}$.
\end{definition}

The last clause is conservative: a newcomer to a partially admitted group is served after the incumbents.

\begin{theorem}[Static misreport]\label{thm:static}
Let $i$ be a price-taking elastic device with sojourn length $L$ and any nondecreasing preference over cumulative volume.
Declaring an RT class $k$ for the whole sojourn instead of $E$ increases its volume by at most $B_{g_i,k}$, i.e.\ by at most $B_{g_i,k}/L$ per slot.
\end{theorem}
\begin{proof}
The deviation yields at most $\sum_t r_{i,k,t}$, and truthful declaration yields at least $\sum_t r_{i,k,t}-B_{g_i,k}$ by Corollary~\ref{cor:window}.
\end{proof}

Without the price-taking hypothesis the same argument gives an identity with an explicit correction: if $R^{\mathrm{dev}}$ is the volume the device receives as a class-$k$ flow on the deviation path and $R^{\mathrm{tr}}=\sum_tr_{i,k,t}$ its reference on the truthful path, then the gain is at most $B_{g_i,k}+(R^{\mathrm{dev}}-R^{\mathrm{tr}})$.
The correction term is the effect of one device on the offers of its tier. We do not bound it analytically; Section~\ref{sec:exp-strategic} measures the total gain exactly, by re-running the system with one device deviating.

\paragraph{Changing class over time.}
A device may also switch its declaration during its stay.

\begin{theorem}[Dynamic declarations]\label{thm:dynamic}
Let $i$ be a price-taking elastic device that declares RT classes during $K$ maximal intervals $J_1,\dots,J_K$ and $E$ otherwise.
The volume it receives during $J_1\cup\dots\cup J_K$ exceeds what truthful declaration yields in the same slots by at most $\sum_{j}B_{g_i,k_j}\le K\max_kB_{g_i,k}$.
If declarations are subject to a hold-down of $\tau$ slots, $K\le\lceil L/(2\tau)\rceil$.
\end{theorem}
\begin{proof}
Apply Corollary~\ref{cor:window} to each window $J_j$ on the truthful path.
\end{proof}

\begin{proposition}[Commitment is necessary]\label{prop:commit}
Consider any allocation rule and a price-taking elastic device that may re-declare its class in every slot.
Its best strategy obtains $\sum_t\max(x_{i,t},r_{i,k,t})$, a gain of $\sum_t(r_{i,k,t}-x_{i,t})^{+}$ over truthful declaration.
This gain is not controlled by the window envy bound: there are paths with class envy at most $1$ over every window on which it grows linearly in $L$.
\end{proposition}
\begin{proof}
The first claim is immediate: declare $k$ exactly in the slots where $r_{i,k,t}>x_{i,t}$.
For the second, let $r_{i,k,t}-x_{i,t}$ alternate between $+1$ and $-1$. Every window sum is at most $1$, and the positive parts add up to $L/2$.
\end{proof}

Proposition~\ref{prop:commit} is not specific to \repay{}.
Any rule that gives RT flows more than their elastic tier-mates in some slots, which is the purpose of having RT classes, can be exploited by a device that is elastic in good slots and ``real-time'' in bad ones.
Deferred compensation and free slot-by-slot re-declaration are incompatible, and some commitment, here the hold-down $\tau$, is the price of offering protection without payments.
Theorem~\ref{thm:dynamic} shows that with commitment the exploitable amount is one envy bound per RT phase.

Theorem~\ref{thm:dynamic} bounds the advantage collected while the device poses as RT. It does not compare the $E$-declared phases, in which the device re-enters with zero deficit while its truthful self would have carried a positive one.
One might expect that less credit can only mean less volume.
This is false: the elastic response \eqref{eq:xresp} is steep near $\lambda-\beta_g=D_i$, so a copy with slightly more credit can be served up to its cap, lose the excess through the reflection in \eqref{eq:D}, and fall behind.
A random search over price-taking environments produced differences of either sign.
We therefore make no analytical claim about these phases and measure the total effect of switching strategies in Section~\ref{sec:exp-strategic}.

\paragraph{The reverse deviation.}
An RT device could declare $E$.
It would lose the guarantee of a rate in $[\rmin_k,\rmax_k]$ in every slot, which is what it values, but it might be served more in some slots.

\begin{proposition}[RT flows do not pose as elastic]\label{prop:reverse}
Assume $u_k'(z)\ge u'(z)$ on $[\rmin_k,\rmax_k]$ (vacuous if $\rmin_k=\rmax_k$) and $u_k(\rmin_k)>\rmin_k\,u'(\rmin_k)$; both hold for the utilities of Section~\ref{sec:model}.
In a slot in which all deficits of tier $g$ are zero, a price-taking class-$k$ flow of tier $g$ obtains at least as much utility from the offer of its group as from the elastic allocation it would receive.
\end{proposition}

The proof is in Appendix~\ref{app:reverse}.
When tier-mates hold positive deficits the RT offer carries a debt price and the statement can fail in a single slot; Section~\ref{sec:exp-strategic} reports the realised gains of this deviation.

\paragraph{Other manipulations.}
Demand caps are measured by the access point, so they are not a strategic variable for downlink traffic.
A subscriber can still inflate its backlog with traffic it does not need, and so accumulate a deficit that it later redeems; since deficits never exceed $B_{g,k}$, so does the credit obtained in this way.
Utilities, deficits and the hold-down are attached to the billing identity, so splitting traffic over several flows of one subscriber changes nothing, and a subscriber is entitled to one RT bundle per class.

\subsection{Near-optimality and the business constraints}\label{sec:opt}

We now show that \repay{} loses little utility relative to the best policy that satisfies the business constraints and is class-envy-free on average.
The statement concerns the fluid allocation, in which a group with admission fraction $a$ earns $n\,a\,u_k(z)$.

\begin{assumption}\label{ass:model}
Let $\xi_t$ denote the \emph{population state} in slot $t$: the active devices with their tiers and declared classes, and the demand caps of the elastic ones.
(i)~The process $(\xi_t)$ is exogenous, i.e.\ not influenced by the allocation, and otherwise arbitrary; at most $N_{\max}$ devices are active at a time.
(ii)~Given $\xi_t=\xi$ and the whole past, $C_t$ is drawn from a distribution $F_\xi$.
(iii)~$|e_{j,t}|\le e_j^{\max}$, demand caps are at most $\ell^{\max}\ge\max_k\rmax_k$, and $0\le U_t\le U^{\max}$.
\end{assumption}

No stationarity or mixing of the population is required; devices may arrive and leave in any pattern.
For a state $\xi$ let $\Pi(\xi)$ be the set of (possibly randomised) maps from $C$ to feasible decisions $(a,y,x)$ of \eqref{eq:slot} such that
\begin{equation}\label{eq:bench}
\E_{F_\xi}[e_j]\le0\ \ \forall j,\qquad
\E_{F_\xi}\big[a_{g_i,k}\,m_{d_i}(z_{g_i,k})-x_i\big]\le0\ \ \forall i\in\cE(\xi),\ k\in\cK,
\end{equation}
and let $U^*(\xi)=\sup_{\pi\in\Pi(\xi)}\E_{F_\xi}[U(\pi)]$.
The second family of constraints is class-envy-freeness \emph{in expectation}: over the capacity fluctuations of a given state, each elastic device receives on average at least the RT offer of its tier.
Policies in $\Pi(\xi)$ may give RT flows priority when capacity is low provided they compensate when it is high, which is precisely the kind of time-shifting that \repay{} performs online.

\begin{theorem}[Utility and constraint violation]\label{thm:opt}
Under Assumption~\ref{ass:model}, if $\Pi(\xi)\neq\emptyset$ for every reachable state, then for all $T\ge1$
\begin{align*}
\frac1T\sum_{t=1}^{T}\E[U_t]\ &\ge\ \frac1T\sum_{t=1}^{T}\E\big[U^*(\xi_t)\big]-\frac{B_0}{V},\\
\frac1T\sum_{t=1}^{T}\E[e_{j,t}]\ &\le\ \sqrt{\frac{2\,(B_0+V\,U^{\max})}{T}}\qquad\text{for each }j,
\end{align*}
where $B_0=\tfrac12\sum_j(e_j^{\max})^2+\tfrac12N_{\max}|\cK|(\ell^{\max})^2$.
\end{theorem}

The proof (Appendix~\ref{app:opt}) uses the Lyapunov function $\tfrac12\sum_jQ_j^2+\tfrac12\sum_{i\in\cE_t}\sum_kD_{i,k}^2$, in which the second sum runs over the devices that are currently present.
Churn is harmless for the drift: a departure removes a nonnegative term and an arrival adds zero.
Together with Theorem~\ref{thm:bound} this gives the familiar trade-off of drift-plus-penalty methods with one difference in kind.
The utility gap is $O(1/V)$ and the violation of the tier constraints vanishes as $O(\sqrt{V/T})$, both in expectation, whereas the envy bound $O(V)$ holds deterministically and for each device over its own, finite, sojourn.

\subsection{What the standard rules do}\label{sec:baselines-theory}

\begin{proposition}\label{prop:baselines}
Consider one tier with one video flow and one saturated elastic device, and constant capacity $C$ with $\rmin_V<C<\min(2\rmin_V,1)$.
\begin{enumerate}[label=(\alph*),nosep]
\item Strict priority for declared RT flows gives the video flow rate $C$ and the elastic device nothing; class envy grows by $C$ per slot, without bound. If the elastic device declares video, only one of the two flows can be served in a slot, so the genuine call is blocked whenever the impostor is served.
\item Class-blind proportional fairness gives both $C/2<\rmin_V$: the call is useless in every slot.
\item Drift-plus-penalty with the same utilities but no deficit queues serves the video flow at rate $C$ and the elastic device never; class envy is unbounded.
\end{enumerate}
\end{proposition}
\begin{proof}
(a) and (b) are direct computations. For (c), if the call is admitted at rate $z$ the objective is $Vw\,[1+\log(z/\rmin_V)+\log(1+C-z)]$, which is increasing in $z$ on $[\rmin_V,C]$ because $C<1$; its value at $z=C$ is at least $Vw$, more than the $Vw\log(1+C)<Vw\log 2$ obtained by blocking the call.
\end{proof}

On the same instance \repay{} serves the call until the elastic device's deficit reaches the threshold of Theorem~\ref{thm:bound}, and from then on gives each side the same long-run volume, which under permanent overload means the call is admitted part of the time.
Bounded class envy and unconditional RT protection cannot both hold when capacity is short forever.
What \repay{} can do, and the three rules above cannot, is protect RT traffic through \emph{transient} shortages at no long-run cost to elastic users.
